\documentclass[UTF8,11pt,a4paper]{ctexart}
\usepackage[margin=24mm]{geometry}
\usepackage{amsmath,amssymb,amsthm,mathtools,bm}
\usepackage{tikz-cd,enumitem,longtable,booktabs}
\usepackage{xurl}
\usepackage[unicode,hidelinks]{hyperref}
\allowdisplaybreaks
\setlength{\emergencystretch}{3em}
\setlength{\parindent}{0pt}
\setlength{\parskip}{0.45em}
\newtheorem*{theorem}{Theorem}
\newtheorem*{lemma}{Lemma}
\newtheorem*{proposition}{Proposition}
\newtheorem*{corollary}{Corollary}
\theoremstyle{definition}\newtheorem*{definition}{Definition}
\theoremstyle{remark}\newtheorem*{remark}{Remark}
\DeclareMathOperator{\Hom}{Hom}
\DeclareMathOperator{\End}{End}
\DeclareMathOperator{\Tr}{Tr}
\DeclareMathOperator{\rank}{rank}
\DeclareMathOperator{\Ind}{Ind}
\DeclareMathOperator{\Res}{Res}
\DeclareMathOperator{\im}{Im}
\DeclareMathOperator{\id}{id}
\newcommand{\R}{\mathbb R}
\newcommand{\C}{\mathbb C}
\newcommand{\Q}{\mathbb Q}
\newcommand{\Z}{\mathbb Z}
\newcommand{\N}{\mathbb N}
\newcommand{\E}{\mathbb E}
\newcommand{\Prob}{\mathbb P}
\newcommand{\eps}{\varepsilon}
\begin{document}
\section*{096｜流形Cech上同调与deRham上同调的比较}
\textbf{作者／集体来源：}Tomasz Mrowka (MIT 18.965 course notes)

\textbf{作品：}Geometry of Manifolds, Fall 2004, Lectures 32--34

\textbf{原文入口：}\url{https://ocw.mit.edu/courses/18-965-geometry-of-manifolds-fall-2004/3fd765fec17898e57e716c7ef400b3f5_lecture34.pdf}

\textbf{原文位置与计题范围：}Lecture 32 §26复形定义；Lecture 33 §27 Lemma 27.1全证明；Lecture 34 §28全比较论证。计题目标为正次数的比较，不将细层无上同调另计一题。

\textbf{获取日期：}2026-09-28。\quad\textbf{复用依据：}MIT OCW CC BY-NC-SA 4.0。

\textbf{整理归属：}下列题面与证明取自所列人类来源，按注明范围转录、摘编；LaTeX环境、标题、换行及中文说明由助手整理。编辑调整在题内说明。

\section*{作者命题与人类证明}

\paragraph{Setting (editorial).}
$M$ is a smooth manifold. We use the source's assumption that countable good covers are cofinal among countable covers, the Poincar\'e lemma, and the long exact cohomology sequence. $\Omega^p$ is the sheaf of smooth $p$-forms; $\mathcal Z^p$ is its subsheaf of closed forms. The counted result is the comparison in positive degrees, obtained from an explicit acyclicity homotopy and dimension shifting.

\subsection*{\v Cech complex (Lecture 32, \S26)}
Let $\mathcal U=\{U_\alpha:\alpha\in A\}$ be an open cover of a topological space. Using the combinatorics of the cover we can define a complex as follows. Let $C^p(\mathcal U)$ be the space of all locally constant functions on $p+1$ fold intersections
\[
U_{\alpha_0}\cap\cdots\cap U_{\alpha_p}
\]
with the symmetry property that if $\sigma$ is a permutation of $0,\ldots,p$, then
\[
f_{\alpha_0\ldots\alpha_p}=\operatorname{sign}(\sigma)f_{\alpha_{\sigma(0)}\ldots\alpha_{\sigma(p)}}.
\]
There is a natural codifferential $\delta:C^p(\mathcal U)\to C^{p+1}(\mathcal U)$ defined by
\[
(\delta f)_{\alpha_0\ldots\alpha_{p+1}}
=\sum_{i=0}^{p+1}(-1)^i
 f_{\alpha_0\ldots\widehat{\alpha_i}\ldots\alpha_{p+1}}
 \big|_{U_{\alpha_0}\cap\cdots\cap U_{\alpha_{p+1}}}.
\]
The direct limit over refinements of the cohomology of these complexes is called the \v Cech cohomology of $M$.

\subsection*{The acyclicity of the sheaf of $p$-forms (Lecture 33, \S27)}
We can consider another version of the \v Cech complex. That is we define $\check C^k(\mathcal U,\Omega^p)$ to be all collections of $p$-forms $\omega_{\alpha_0\ldots\alpha_k}$ defined on $U_{\alpha_0\ldots\alpha_k}$ with the symmetry properties above. The same formula defines a differential mapping
\[
\check C^k(\mathcal U,\Omega^p)\longrightarrow\check C^{k+1}(\mathcal U,\Omega^p).
\]
\begin{lemma}[27.1]
Given an open cover $\mathcal U$, the \v Cech complex
\[
\cdots\longrightarrow\check C^{k-1}(\mathcal U;\Omega^p)
\xrightarrow{\delta}\check C^k(\mathcal U;\Omega^p)
\xrightarrow{\delta}\check C^{k+1}(\mathcal U;\Omega^p)\longrightarrow\cdots
\]
is exact so long as $k>0$.
\end{lemma}
\begin{proof}
Fix a partition of unity $\{\phi_\beta:\beta\in B\}$ subordinate to $\mathcal U=\{U_\alpha\}_{\alpha\in A}$. The supports of the $\phi_\beta$ are a refinement of the $U_\alpha$ and we choose a refinement function $r:B\to A$ so that $\operatorname{supp}(\phi_\beta)\subset U_{r(\beta)}$. Define
\[
K:\check C^k(\mathcal U;\Omega^p)\longrightarrow\check C^{k-1}(\mathcal U;\Omega^p)
\]
by
\[
(K\omega)_{\alpha_0\ldots\alpha_{k-1}}
=\sum_{\beta\in B}\phi_\beta\omega_{r(\beta)\alpha_0\ldots\alpha_{k-1}}.
\]
Since the supports of the $\phi_\beta$'s are locally finite by definition of partition of unity, this is well defined. Now consider, where $k\ge1$,
\begin{align*}
((\delta K+K\delta)\omega)_{\alpha_0\ldots\alpha_k}
&=\sum_{i=0}^k(-1)^i(K\omega)_{\alpha_0\ldots\widehat{\alpha_i}\ldots\alpha_k}
 +\sum_{\beta\in B}\phi_\beta(\delta\omega)_{r(\beta)\alpha_0\ldots\alpha_k}\\
&=\sum_{i=0}^k\sum_{\beta\in B}(-1)^i\phi_\beta
\omega_{r(\beta)\alpha_0\ldots\widehat{\alpha_i}\ldots\alpha_k}\\
&\quad+\sum_{\beta\in B}\phi_\beta\omega_{\alpha_0\ldots\alpha_k}
-\sum_{\beta\in B}\sum_{j=0}^k(-1)^j\phi_\beta
\omega_{r(\beta)\alpha_0\ldots\widehat{\alpha_j}\ldots\alpha_k}\\
&=\omega_{\alpha_0\ldots\alpha_k}.
\end{align*}
We have used that the sum is locally finite to rearrange the order of summation. Thus we have proved the identity is cochain homotopic to zero and so the cohomology groups are zero. Note that if $k=0$ the argument proves nothing.
\end{proof}

\subsection*{The comparison (Lecture 34, \S28)}
\begin{theorem}[\v Cech--de Rham comparison]
There is a natural isomorphism
\[
\check H^p(M;\R)\simeq H^p_{\mathrm{deR}}(M;\R),\qquad p\ge1.
\]
\end{theorem}
\begin{proof}
The scheme of the proof is to first restrict attention to countable good covers which we assume to be cofinal in the set of countable covers. The Poincar\'e lemma tells us that for a contractible open set $U$
\[
\R\hookrightarrow\Omega^0(U)\xrightarrow d\Omega^1(U)
\xrightarrow d\Omega^2(U)\xrightarrow d\cdots
\]
is a long exact sequence. We introduce the notation $\mathcal Z^p$ for the closed $p$-forms so that
\[
\mathcal Z^p(U)=\{\theta\in\Omega^p(U):d\theta=0\}.
\]
Then we can break up this long exact sequence into short exact sequences:
\[
0\longrightarrow\mathcal Z^p(U)\hookrightarrow\Omega^p(U)
\xrightarrow d\mathcal Z^{p+1}(U)\longrightarrow0.
\]
Note that $\mathcal Z^0(U)$ is the constant functions, so a copy of $\R$. These induce long exact sequences in cohomology,
\begin{multline*}
\check H^{i-1}(M;\Omega^p)\longrightarrow\check H^{i-1}(M;\mathcal Z^{p+1})
\longrightarrow\check H^i(M;\mathcal Z^p)\\
\longrightarrow\check H^i(M;\Omega^p)\longrightarrow\cdots.
\end{multline*}
We have seen that $\check H^i(M;\Omega^p)=\{0\}$ for $i>0$, and hence
\[
\check H^i(M;\mathcal Z^p)\simeq\check H^{i-1}(M;\mathcal Z^{p+1})
\]
for $i\ge2$. Now by definition the $p$th \v Cech cohomology group of $M$ is
\[
\check H^p(M;\R)=\check H^p(M;\mathcal Z^0).
\]
Repeatedly applying the isomorphism above, we have
\[
\check H^p(M;\R)\simeq\check H^1(M;\mathcal Z^{p-1}).
\]
Now consider the beginning of the long exact sequence,
\begin{multline*}
0\longrightarrow H^0(M;\mathcal Z^{p-1})\longrightarrow H^0(M;\Omega^{p-1})
\longrightarrow H^0(M;\mathcal Z^p)\\
\longrightarrow H^1(M;\mathcal Z^{p-1})\longrightarrow H^1(M;\Omega^{p-1}),
\end{multline*}
which becomes
\[
0\longrightarrow\mathcal Z^{p-1}(M)\longrightarrow\Omega^{p-1}(M)
\xrightarrow d\mathcal Z^p(M)\longrightarrow H^1(M;\mathcal Z^{p-1})\longrightarrow0.
\]
Thus
\[
H^1(M;\mathcal Z^{p-1})\simeq\mathcal Z^p(M)/d\Omega^{p-1}(M)
=H^p_{\mathrm{deR}}(M;\R).
\]
Thus we have proved that there is a natural isomorphism
$\check H^p(M;\R)\simeq H^p_{\mathrm{deR}}(M;\R)$.
\end{proof}

\section*{整理说明与可修改标注（助手）}
\textbf{标准先修：}光滑流形与单位分解；Poincare引理；良覆盖共尾性；层短正合列引出的上同调长正合列。这些是作者比较证明明确采用的前提。

\textbf{本题仍需完成的工作：}从局部单位分解写出Cech复形上的降次数同伦，完整消去高阶形式层上同调，再把闭形式的正合列逐次移位为deRham商空间。

\textbf{取材说明：}不收入Lecture 32中未写完的Poincare证明，而把该引理明确列为先修。转录统一了K的次数标记与其分量公式（降一次），并把Lecture 34倒数第二条长正合列末项原印的次数p改为与该列一致的次数p减一；均为就地注明的下标整理。其余证明与作者顺序一致。讲义33第1页和34已查看原页图像；33第2页图像请求失败，已读取其完整文字层。

\textbf{$c$：}流形上两种上同调的比较

\textbf{$a$：}单位分解；Cech复形；链同伦；微分形式；Poincare引理；长正合列与维数移位

\texttt{cross\_theory\_status = yes}

\textbf{具体联系及位置：}Lecture 33的K将光滑单位分解作用到覆盖的组合复形；Lecture 34再利用局部微分正合性与全局上同调连接映射，识别两种不同表示的拓扑不变量。

这些标注是非穷尽初稿；未列出不表示未使用。
\end{document}
